\documentclass[11pt,a4paper]{article}

\usepackage[margin=2.7cm]{geometry}
\usepackage[T1]{fontenc}
\usepackage{lmodern}
\usepackage{amsmath,amssymb,amsthm}
\usepackage{bbm}
\usepackage{graphicx}
\usepackage{booktabs}
\usepackage{array}
\usepackage[dvipsnames]{xcolor}
\usepackage[colorlinks=true,linkcolor=MidnightBlue,citecolor=MidnightBlue,urlcolor=MidnightBlue]{hyperref}
% TeX Live 2025's microtype 3.2a warns `Command \showhyphens has changed'
% against this kernel even in an empty document.  Mute the kernel warning
% during microtype's begin-document check only (hooks run in registration
% order: mute -> microtype's check -> restore).
\makeatletter
\AddToHook{begindocument/before}{\let\ma@saved@warning\@latex@warning@no@line
  \let\@latex@warning@no@line\@gobble}
\makeatother
\usepackage{microtype}
\makeatletter
\AddToHook{begindocument}{\let\@latex@warning@no@line\ma@saved@warning}
\makeatother
\emergencystretch=1.5em

\newtheorem{proposition}{Proposition}
\newtheorem{lemma}{Lemma}
\theoremstyle{definition}
\newtheorem{definition}{Definition}
\newtheorem{remark}{Remark}

% units
\newcommand{\hPa}{\,\mathrm{hPa}}
\newcommand{\Jkg}{\,\mathrm{J\,kg^{-1}}}
\newcommand{\JkgK}{\,\mathrm{J\,kg^{-1}\,K^{-1}}}
\newcommand{\K}{\,\mathrm{K}}
\newcommand{\ms}{\,\mathrm{m\,s^{-1}}}
% thermodynamic symbols
\newcommand{\Rd}{R_d}
\newcommand{\Rv}{R_v}
\newcommand{\cpd}{c_{pd}}
\newcommand{\cpv}{c_{pv}}
\newcommand{\cl}{c_\ell}
\newcommand{\Lv}{L_v}
\newcommand{\rt}{r_t}
\newcommand{\rvv}{r_v}
\newcommand{\rl}{r_\ell}
\newcommand{\es}{e_s}
\newcommand{\pd}{p_d}
\newcommand{\eps}{\varepsilon}
\newcommand{\dd}{\mathrm{d}}
\newcommand{\Trho}{T_\rho}

\title{Thermodynamic foundations of the moist-parcel entropy in \texttt{tcpyPI}\\
\large A first-principles derivation of \texttt{entropy\_S}, its conservation,\\
and the temperature inversion on the reversible adiabat}
\author{Companion to \emph{An artificial discontinuity in the tcpyPI pressure solver}}
\date{July 14, 2026}

\begin{document}
\maketitle

\begin{abstract}
\noindent
The potential-intensity solver in \texttt{tcpyPI} lifts an air parcel along a
\emph{reversible moist adiabat} and, at each pressure level, recovers the
parcel temperature by holding one quantity fixed: the total specific entropy
computed by \texttt{utilities.entropy\_S}. This note derives that quantity, and
everything around it, from first principles. We start from the statistical
definition of entropy (Boltzmann/Gibbs), pass through the Sackur--Tetrode
equation to the ideal-gas form $s=c_p\ln T-R\ln p+\mathrm{const}$, treat the
parcel as a multicomponent Dalton mixture with a condensed phase, derive the
Clausius--Clapeyron relation and the latent-heat entropy from equality of
chemical potentials, and assemble Emanuel's total specific entropy
(E94, eq.~4.5.9)---the exact expression in the code. We then prove it is
conserved on the ascent (adiabatic $\Rightarrow$ no entropy flux; reversible
$\Rightarrow$ no entropy production), give the microscopic Liouville/adiabatic
-theorem reading of that conservation, and finish with the moist adiabat and
the well-posed \emph{inversion} $s(T,p)=s_0\mapsto T$ that the solver performs
by Newton iteration. Every formula is tied back to the specific function and
constant it corresponds to in the source. Numerical constants are quoted as in
\texttt{tcpyPI/constants.py}; SI throughout, pressures excepted (hPa).
\end{abstract}

\tableofcontents

\section{Scope and the target formula}\label{sec:scope}

The object we are after is the total specific entropy per unit mass of
\emph{dry} air, as coded in \texttt{utilities.entropy\_S(T,R,P)}:
\begin{equation}\label{eq:target}
s \;=\; (\cpd+\rt\,\cl)\,\ln T \;-\; \Rd\,\ln(p-e) \;+\; \frac{\Lv\,\rvv}{T}
     \;-\; \rvv\,\Rv\,\ln H,
\end{equation}
with $T$ the temperature (K), $p$ the total pressure, $e$ the partial
pressure of water vapour, $p-e=\pd$ the partial pressure of dry air,
$\rt$ the total-water mixing ratio, $\rvv$ the vapour mixing ratio, $\Lv$ the
latent heat of vaporisation, $H=e/\es$ the relative humidity (capped at $1$),
and $\cpd,\cl,\Rd,\Rv$ material constants (Table~\ref{tab:const}). The code
evaluates \eqref{eq:target} at the parcel's launch state, where the air is
unsaturated so $\rvv=\rt=R$; that is why a single argument \texttt{R} feeds all
water terms. When the parcel later saturates, the same $s$ is held fixed and
$T$ is solved for---the \emph{inversion} of \S\ref{sec:inversion}.

Our task is to show where \eqref{eq:target} comes from, why it is the right
thing to conserve, and why the inversion is well posed. We build it in five
movements: statistical mechanics (\S\ref{sec:statmech}), classical
thermodynamic scaffolding (\S\ref{sec:classical}), the multicomponent parcel
(\S\ref{sec:parcel}), phase equilibrium and latent heat (\S\ref{sec:phase}),
and the assembly (\S\ref{sec:assembly}). Conservation
(\S\ref{sec:conservation}), the moist adiabat (\S\ref{sec:adiabat}) and the
inversion (\S\ref{sec:inversion}) follow.

\begin{table}[htb]
\centering
\small
\begin{tabular}{llr}
\toprule
Symbol & Meaning & Value (\texttt{constants.py})\\
\midrule
$\cpd$ & specific heat of dry air at const.\ $p$ & $1005.7\JkgK$\\
$\cpv$ & specific heat of water vapour at const.\ $p$ & $1870.0\JkgK$\\
$\cl$  & (modified) specific heat of liquid water & $2500.0\JkgK$\\
$\Rd$  & gas constant of dry air & $287.04\JkgK$\\
$\Rv$  & gas constant of water vapour & $461.5\JkgK$\\
$\eps=\Rd/\Rv$ & mass ratio of the two gases & $0.6220$\\
$\Lv^{0}$ & latent heat at $0^{\circ}$C & $2.501\times10^{6}\Jkg$\\
$\cpv-\cl$ & Kirchhoff slope $\dd\Lv/\dd T$ & $-630\JkgK$\\
\bottomrule
\end{tabular}
\caption{Material constants, quoted from \texttt{tcpyPI/constants.py}. The
``modified'' $\cl=2500$ (rather than $4190$) is a tuning of the reference,
discussed in Remark~\ref{rem:modcl}.}
\label{tab:const}
\end{table}

\section{The statistical origin of entropy}\label{sec:statmech}

\subsection{The operator picture underneath}\label{ssec:operator}

Before the classical development, it is worth stating the exact
quantum-statistical object that development approximates, because it removes
several apparent ambiguities in the symbol ``$\sum_i$'' that a careful reader
will otherwise trip on. Rigorously, energy is not a number but the Hamiltonian
\emph{operator} $\hat H$ acting on the $N$-particle Hilbert space
$\mathcal H_N$, and the partition function is a trace,
\begin{equation}\label{eq:trace}
Z=\operatorname{Tr}e^{-\beta\hat H},
\end{equation}
whose value is independent of the basis in which it is evaluated. The microstate
energies $\{E_i\}$ are the eigenvalues of $\hat H$; the index $i$ labels an
orthonormal energy eigenbasis; and the classical sum $\sum_i e^{-\beta E_i}$
used throughout this section is exactly \eqref{eq:trace} written in that
eigenbasis, where $\hat H$ is diagonal with entries $E_i$. Three things this
clarifies at once:
\begin{itemize}
\item \textbf{The index $i$ is discrete, and $\mathcal H_N,\hat H$ depend on
$(V,N)$.} Confinement to a finite volume discretises the spectrum---a particle
in a box of side $V^{1/3}$ has $\varepsilon\propto(n_x^2+n_y^2+n_z^2)\,V^{-2/3}$
with $n\in\mathbb Z_+^3$---so $i$ is a genuine countable label, not a continuum.
The parameters $V$ and $N$ do \emph{not} index microstates; they select which
Hilbert space and Hamiltonian one is on. Fixing $(V,N)$, as the canonical
ensemble does, means tracing over the eigenbasis of one definite operator
$\hat H_{V,N}$.
\item \textbf{The classical sum is the $\hbar\to0$ limit.} The trace
\eqref{eq:trace} reduces to the phase-space integral
$(N!\,h^{3N})^{-1}\!\int\dd\Gamma\,e^{-\beta H_{\mathrm{cl}}}$ of
\S\ref{ssec:lambda}, with relative corrections $O(\hbar^2)$. The factors
$h^{3N}$ (one quantum state per phase-space cell of volume $h^{3N}$) and $N!$
(indistinguishability, from the (anti)symmetrisation
$\mathcal H_N\subset\mathcal H_1^{\otimes N}$) are \emph{not} classical
inputs---they are residues of the quantum count, which is why they enter
\eqref{eq:ztrans2}--\eqref{eq:ZN} seemingly by fiat.
\item \textbf{Everything is a spectral functional of $\hat H$.} $Z(\beta)$ is
the Laplace transform of the density of states
$\rho(E)=\operatorname{Tr}\delta(E-\hat H)$, and $F,\langle E\rangle,S$ all
follow from it by one derivative each---spelled out in \S\ref{ssec:Z},
\[
F=-k_BT\ln Z,\qquad
\langle E\rangle=-\frac{\partial\ln Z}{\partial\beta},\qquad
S=\frac{\langle E\rangle-F}{T}=-\frac{\partial F}{\partial T}
\]
---so thermodynamics is the spectral geometry of the Hamiltonian. For
$N$ free identical particles the canonical $Z_N$ is, exactly, a symmetric
function of the single-particle Boltzmann weights $x_k=e^{-\beta\varepsilon_k}$
(one weight per single-particle level $k$), and \emph{which} symmetric function
records the exchange statistics:
\begin{itemize}
\item the elementary $e_N(x)=\sum_{k_1<\dots<k_N}x_{k_1}\cdots x_{k_N}$ for
\textbf{fermions} (strictly increasing indices $\Rightarrow$ no level used
twice---Pauli exclusion);
\item the complete homogeneous
$h_N(x)=\sum_{k_1\le\dots\le k_N}x_{k_1}\cdots x_{k_N}$ for \textbf{bosons}
($\le\Rightarrow$ any level may repeat, any occupation allowed);
\item the leading power-sum term $p_1(x)^N/N!$, with $p_1(x)=\sum_k x_k=z$, for
the \textbf{Maxwell--Boltzmann count}---the classical $z^N/N!$ of
\eqref{eq:ZN} (derived in \S\ref{ssec:Z}): particles taken as independent and
then divided by $N!$ for indistinguishability, the dilute, non-degenerate limit
(\S\ref{ssec:lambda}) in which the fermion/boson distinction washes out.
\end{itemize}
So the indistinguishability that forces the $N!$ is the $S_N$-symmetry of these
functions.
\end{itemize}
Finally, the various \emph{ensembles} are Legendre transforms of one another,
each holding fixed a different member of the conjugate pairs $(E,T)$, $(V,p)$,
and $(N,\mu)$---energy with temperature, volume with pressure $p$, and particle
number with chemical potential $\mu$ (the free-energy cost of adding one
particle); ``summing over'' a mechanical variable trades it for its conjugate:
\begin{center}
\small
\begin{tabular}{lll}
\toprule
sum runs over & held fixed & potential\\
\midrule
microstates ($E$ free) & $E,V,N$ & $S(E,V,N)$\quad(microcanonical)\\
microstates & $T,V,N$ & $F(T,V,N)$\quad(canonical)\\
microstates $+\;N$ & $T,V,\mu$ & $\Omega(T,V,\mu)$\quad(grand)\\
microstates $+\;V$ & $T,p,N$ & $G(T,p,N)$\quad(isobaric)\\
\bottomrule
\end{tabular}
\end{center}
The canonical development below fixes $(T,V,N)$ and works with $F$; the parcel
of the companion analysis ultimately wants $(T,p)$---the Gibbs potential
$G$---which is why the specific entropy is finally expressed as $s(T,p)$ and
inverted for $T$ at prescribed $p$. What follows carries out \eqref{eq:trace} in
its classical $\hbar\to0$ form; the operator identity is the exact statement
being approximated.

\subsection{Boltzmann and Gibbs, and the system we will use}\label{ssec:BG}

Entropy is, microscopically, a count of microstates. Fix the system once, for
the whole section: a gas of $N$ particles, each of mass $m$, confined to a
container of volume $V$ and held in equilibrium at temperature $T$. Its
\emph{macrostate} is specified by the three numbers $(T,V,N)$; a
\emph{microstate} $i$ is a complete specification of every particle's position
and momentum (classically) or of the joint quantum state (quantally), and has a
definite total energy $E_i$. This distinction---$E_i$ is the energy of the
\emph{whole system} in microstate $i$, whereas $\varepsilon$ (introduced in
\S\ref{ssec:lambda}) will denote the energy of a \emph{single particle}---is
used consistently below.

For an \emph{isolated} system all $\Omega$ microstates consistent with the
macrostate are equally probable (the fundamental postulate of statistical
mechanics), and Boltzmann's principle is
\begin{equation}\label{eq:boltzmann}
S \;=\; k_B \ln \Omega ,\qquad k_B=1.381\times10^{-23}\,\mathrm{J\,K^{-1}} .
\end{equation}
A parcel is \emph{not} isolated---it exchanges energy with its surroundings---so
its microstates are not equiprobable; they carry probabilities $p_i$, and the
entropy is the more general Gibbs (information) form
\begin{equation}\label{eq:gibbs}
S \;=\; -k_B \sum_i p_i \ln p_i ,
\end{equation}
which reduces to \eqref{eq:boltzmann} when the $p_i$ are uniform, $p_i=1/\Omega$.
To use \eqref{eq:gibbs} we must first find the $p_i$.

\subsection{The Boltzmann distribution from energy conservation}\label{ssec:boltzdist}

The $p_i$ are fixed by one physical input---\emph{conservation of energy}
between the system and its surroundings---plus the postulate above. Put the
system in weak thermal contact with a large reservoir (subscript $\mathrm{res}$
throughout this subsection; not to be confused with the specific gas constant
$R$ that appears bare from \S\ref{ssec:intensive} onward), and treat the
composite (system $+$ reservoir) as isolated with a fixed total energy
\begin{equation}\label{eq:econs}
E_{\mathrm{tot}}=E_i+E_{\mathrm{res}}=\mathrm{const}.
\end{equation}
``Weak contact'' means the interaction energy is negligible, so the energies are
additive \eqref{eq:econs}; ``isolated composite'' lets us apply the equal-\emph{a
priori}-probability postulate to it. When the system sits in one definite
microstate $i$ (energy $E_i$), the reservoir may be in any of its
$\Omega_{\mathrm{res}}(E_{\mathrm{tot}}-E_i)$ microstates, all equally likely. Hence the
probability of system-microstate $i$ is proportional to that reservoir count,
\begin{equation}\label{eq:pprop}
p_i \;\propto\; \Omega_{\mathrm{res}}\!\big(E_{\mathrm{tot}}-E_i\big).
\end{equation}
Because the reservoir is large, $E_i\ll E_{\mathrm{tot}}$, so expand its
\emph{logarithm} (the extensive quantity) to first order:
\begin{equation}\label{eq:expand}
\ln\Omega_{\mathrm{res}}(E_{\mathrm{tot}}-E_i)
\;\simeq\;\ln\Omega_{\mathrm{res}}(E_{\mathrm{tot}})
-E_i\underbrace{\frac{\partial\ln\Omega_{\mathrm{res}}}{\partial E}\Big|_{E_{\mathrm{tot}}}}_{\displaystyle\equiv\,\beta}.
\end{equation}
The coefficient $\beta$ is a property of the reservoir alone; by
\eqref{eq:boltzmann} applied to the reservoir, $S_{\mathrm{res}}=k_B\ln\Omega_{\mathrm{res}}$, and the
thermodynamic definition of temperature $1/T=\partial S_R/\partial E$, it is
\begin{equation}\label{eq:beta}
\beta=\frac{1}{k_B}\frac{\partial S_{\mathrm{res}}}{\partial E}=\frac{1}{k_BT}.
\end{equation}
Exponentiating \eqref{eq:expand} and absorbing the $E_i$-independent factor into
the normalisation gives the \emph{Boltzmann distribution}
\begin{equation}\label{eq:boltzdist}
p_i=\frac{e^{-\beta E_i}}{Z},\qquad
Z=\sum_i e^{-\beta E_i},
\end{equation}
with the \emph{partition function} $Z$ enforcing $\sum_i p_i=1$. The
assumptions are now explicit: (i) equal \emph{a priori} probabilities for the
isolated composite; (ii) conservation of energy \eqref{eq:econs}; (iii) a
reservoir large enough that the linear term in \eqref{eq:expand} suffices; (iv)
weak coupling, so energies add. (Equivalently, by Jaynes' maximum-entropy
argument, \eqref{eq:boltzdist} is the distribution that maximises
\eqref{eq:gibbs} subject only to $\sum_i p_i=1$ and a fixed mean energy
$\langle E\rangle=\sum_i p_iE_i$, with $\beta$ the Lagrange multiplier for
energy---the two derivations agree.)

\subsection{The reservoir from the inside: canonical typicality and the
\texorpdfstring{$n$}{n}-oscillator example}\label{ssec:typicality}

The derivation just given works, but three of its four inputs are uncomfortable
to a reader who believes the world is one closed system evolving unitarily. The
reservoir is an external object introduced by hand; the equal-\emph{a priori}
-probability postulate~(i) is imposed, not derived; and the ensemble language
treats $p_i$ as a subjective distribution over copies rather than a property of
a single physical state. There is a modern reading---\emph{canonical
typicality} (Popescu, Short and Winter, \emph{Nat.\ Phys.}~\textbf{2}, 754
(2006); Goldstein, Lebowitz, Tumulka and Zangh\`i, \emph{Phys.\ Rev.\ Lett.}
\textbf{96}, 050403 (2006))---that produces the \emph{same} Boltzmann weights
from a single pure state of one closed system, with the postulate replaced by a
theorem. It is worth a page, both because it dissolves the discomfort and
because the oscillator makes every object in it explicit.

\paragraph{The general statement.} Split the closed system into the piece we
observe and the rest, $\mathcal H=\mathcal H_S\otimes\mathcal H_E$ (``$S$'' the
system, ``$E$'' its environment; dimensions $d_S,d_E$), and impose a global
constraint---here a fixed total energy in a thin shell---by restricting to a
subspace $\mathcal H_R\subseteq\mathcal H$ of dimension $d_R=\dim\mathcal H_R$.
This $d_R$ is exactly the microstate count $\Omega$ of \eqref{eq:boltzmann}, now
read as a Hilbert-space dimension. The uniform (microcanonical) mixture on the
shell is $\rho_{\mathrm{mc}}=\mathbbm{1}_R/d_R$, and its reduction defines what PSW
call the \emph{canonical state} of the system,
\begin{equation}\label{eq:psw-can}
\varrho_S \;\equiv\; \operatorname{Tr}_E\,\rho_{\mathrm{mc}}
        \;=\; \operatorname{Tr}_E\frac{\mathbbm{1}_R}{d_R}.
\end{equation}
Here $\operatorname{Tr}_E$ is the \emph{partial} trace over the environment---the
linear map from operators on $\mathcal H_S\otimes\mathcal H_E$ to operators on
$\mathcal H_S$ fixed on product operators by
\begin{equation}\label{eq:ptrace}
\operatorname{Tr}_E(A\otimes B)=A\,\operatorname{Tr}B
\qquad\Big(\text{contrast the full trace }
\operatorname{Tr}(A\otimes B)=\operatorname{Tr}A\,\operatorname{Tr}B\Big),
\end{equation}
and on a general operator by linearity. It contracts the $E$ factor to a number
while leaving the $S$ factor intact, so---unlike the full trace---it returns an
\emph{operator} on $\mathcal H_S$, not a scalar; \eqref{eq:psw-can} is thus a
genuine density matrix on $S$, as is the pure-state reduction introduced next.
Now take a \emph{single} pure state $|\phi\rangle\in\mathcal H_R$---one definite
vector, no ensemble---and form its reduced density matrix
$\rho_S=\operatorname{Tr}_E|\phi\rangle\langle\phi|$. The theorem is that for the
overwhelming majority of $|\phi\rangle$ (Haar measure on the unit sphere of
$\mathcal H_R$), $\rho_S$ is indistinguishable from \eqref{eq:psw-can}:
\begin{equation}\label{eq:psw-bound}
\big\langle\,\lVert\rho_S-\varrho_S\rVert_1\,\big\rangle_{\phi}
 \;\le\; \sqrt{\frac{d_S}{d_E^{\mathrm{eff}}}},
\qquad
d_E^{\mathrm{eff}}\equiv\frac{1}{\operatorname{Tr}(\varrho_E^{2})},\quad
\varrho_E=\operatorname{Tr}_S\rho_{\mathrm{mc}},
\end{equation}
and, by L\'evy's lemma (concentration of measure on the high-dimensional
sphere), the fraction of states violating this by any fixed margin is
exponentially small in $d_R$ (nonempty but exponentially rare---the exceptional
states, and the SHO's starkest one, are pinned down in
Remark~\ref{rem:atypical}). The equal-\emph{a priori} postulate is thereby
\emph{demoted to a theorem}: one need not assume a flat distribution over
microstates---almost every individual pure state already \emph{looks} flat to a
small subsystem, because that is what generic entanglement across a large
bipartition does. The weights $p_i$ become eigenvalues of a reduced density
matrix (an objective property of one state), not a subjective ensemble.

\paragraph{The $n$-oscillator realisation.} Make
$\mathcal H_S,\mathcal H_E,\mathcal H_R$ concrete with the harmonic oscillator,
the one closed system where the counting is elementary. Take $n$ identical
oscillators
\begin{equation}\label{eq:sho-H}
\hat H=\hbar\omega\sum_{i=1}^{n}\hat n_i\qquad(\text{zero-point dropped}),
\end{equation}
equivalently one particle in an isotropic $n$-dimensional harmonic trap. Let $S$
be oscillator~$1$ and $E$ the remaining $n-1$; the global constraint is a fixed
total number of quanta $M=\sum_i k_i$, i.e.\ total energy
$E_{\mathrm{tot}}=\hbar\omega M$. The shell
$\mathcal H_R=\mathrm{span}\{\,|k_1,\dots,k_n\rangle:\sum_i k_i=M\,\}$ has the
stars-and-bars dimension
\begin{equation}\label{eq:sho-dR}
d_R=\Omega(M,n)=\binom{M+n-1}{\,n-1\,},
\end{equation}
the microcanonical count for this system. A generic normalised state in it is
$|\phi\rangle=\sum_{\sum k_i=M}c_{k_1\cdots k_n}\,|k_1,\dots,k_n\rangle$ with
Haar-random coefficients.

\emph{The reduced density matrix is automatically diagonal.} Writing an
environment configuration as $\vec K=(k_2,\dots,k_n)$ with
$|\vec K|\equiv\sum_{i\ge2}k_i$, the constraint forces $k_1=M-|\vec K|$ for
\emph{every} term, so each $\vec K$ pairs with a single system level. Hence
\begin{equation}\label{eq:sho-rhoS}
\rho_S=\operatorname{Tr}_E|\phi\rangle\langle\phi|
      =\sum_{k=0}^{M}\Big(\,\underbrace{\textstyle\sum_{|\vec K|=M-k}
        |c_{k,\vec K}|^{2}}_{\displaystyle \rho_S(k)}\Big)\,|k\rangle\langle k|,
\end{equation}
with \emph{no} off-diagonal entries: two different system energies would need
the same environment state to carry two different energies, which the global
conservation law forbids. This is dephasing by a conservation law made
concrete---the microscopic reason $\rho_S$ is a classical distribution over
energy levels rather than a coherent superposition, so that the $\{p_i\}$ of
\eqref{eq:boltzdist} are literally its diagonal.

\emph{Its average is the canonical state.} For a Haar-random unit vector in
$d_R$ complex dimensions $\langle|c_{k,\vec K}|^{2}\rangle=1/d_R$, and the number
of environment configurations with $|\vec K|=M-k$ is $\Omega(M-k,\,n-1)$, so
\begin{equation}\label{eq:sho-avg}
\big\langle\rho_S(k)\big\rangle
 =\frac{\Omega(M-k,\,n-1)}{d_R}
 =\frac{\dbinom{M-k+n-2}{\,n-2\,}}{\dbinom{M+n-1}{\,n-1\,}}
 \;\equiv\;\varrho_S(k),
\end{equation}
which is exactly \eqref{eq:psw-can} written out---the marginal of the
microcanonical shell onto one oscillator.

\emph{And a single state hugs that average.} The diagonal entry $\rho_S(k)$ is a
sum of $N_k=\Omega(M-k,n-1)$ terms $|c|^{2}$, each $\sim1/d_R$; for a random unit
vector its relative fluctuation is $1/\sqrt{N_k}$, exponentially small in $n$.
The PSW bound \eqref{eq:psw-bound} makes this sharp. Here the reduced
microcanonical environment state $\varrho_E$ is uniform over the $d_R$
environment configurations with $|\vec K|\le M$ (hockey-stick identity
$\sum_{j=0}^{M}\Omega(j,n-1)=d_R$), so $\operatorname{Tr}(\varrho_E^{2})=1/d_R$
and $d_E^{\mathrm{eff}}=d_R$. With $d_S=M+1$ accessible levels,
\begin{equation}\label{eq:sho-bound}
\big\langle\lVert\rho_S-\varrho_S\rVert_1\big\rangle
 \le\sqrt{\frac{M+1}{\dbinom{M+n-1}{\,n-1\,}}}
 \;\approx\;\sqrt{\mu n}\;e^{-n\sigma(\mu)/2},
\qquad
\sigma(\mu)=(1{+}\mu)\ln(1{+}\mu)-\mu\ln\mu,
\end{equation}
where $\mu=M/n$ is the mean occupancy and $\sigma(\mu)$ is---not
coincidentally---the canonical entropy per oscillator in units of $k_B$. The
exponent is $n$ times a thermodynamic entropy: the measure of atypical states is
$e^{-S/k_B}$. Numbers make the point brutal. For a modest bath of $n=100$
oscillators at $\mu=1$ ($\beta\hbar\omega=\ln2$),
$d_R=\binom{199}{99}\approx4.5\times10^{58}$ and the bound is below $10^{-28}$: a
\emph{single} random pure state of one hundred oscillators already presents
oscillator~$1$ with a density matrix that is thermal to twenty-eight decimal
places---no reservoir, no ensemble, no postulate, only the geometry of a
high-dimensional sphere.

\begin{remark}[where canonical typicality fails]\label{rem:atypical}
The bound \eqref{eq:psw-bound}/\eqref{eq:sho-bound} governs the \emph{typical}
state; its exceptional set is nonempty, and the oscillator names it. Typicality
fails precisely for the states with little $S{:}E$ entanglement: a product
$|\psi_S\rangle\otimes|\psi_E\rangle$ has $\rho_S=|\psi_S\rangle\langle\psi_S|$
pure, as far from thermal as a state can be. The extreme case is a single shell
Fock vector with every quantum in the system,
\begin{equation}\label{eq:atypical}
|\phi\rangle=|M,0,\dots,0\rangle,\qquad
\rho_S=|M\rangle\langle M|,\qquad
\tfrac12\lVert\rho_S-\varrho_S\rVert_1=1-P(M)=1-\tfrac1{d_R},
\end{equation}
using $P(M)=\Omega(0,n-1)/d_R=1/d_R$. The trace distance is $\approx1$, the
largest two density matrices can have, whereas the typical value
\eqref{eq:sho-bound} is $\sim10^{-28}$---this lone state overshoots the bound by
some twenty-eight orders of magnitude. Yet it costs the theorem nothing: such
product states form a submanifold far below the sphere's $2d_R-1$ real
dimensions and are met with probability $\sim e^{-cn}$. They are exactly the
definite-per-oscillator-occupation states---the natural stationary states of the
\emph{integrable} oscillator---so this kinematic exception is the shadow of the
dynamical non-thermalisation in the cautionary tale below, and the reason a
chaotic (ETH-obeying) $\hat H$, whose eigenstates are \emph{themselves} typical,
has no such exceptions.
\end{remark}

\paragraph{Recovering the Boltzmann weight, and meeting \S\ref{ssec:boltzdist}.}
It remains to identify $\varrho_S$ with the Gibbs state. This is the
\emph{equivalence of ensembles}, and for the oscillator it is one line: the
consecutive ratio of \eqref{eq:sho-avg} is
\begin{equation}\label{eq:sho-ratio}
\frac{\varrho_S(k{+}1)}{\varrho_S(k)}=\frac{M-k}{\,M-k+n-2\,}
\xrightarrow[\;M=\mu n\;]{n\to\infty}\frac{\mu}{\mu+1}\equiv x,
\end{equation}
so $\varrho_S(k)\to(1-x)\,x^{k}$, a geometric (Planck) distribution---precisely
$e^{-\beta\hbar\omega\,k}/Z_S$ with $Z_S=(1-x)^{-1}$ and $x=e^{-\beta\hbar\omega}$.
The temperature is not an input: differentiating \eqref{eq:sho-dR} gives
$\beta=\partial\ln\Omega/\partial E=(\hbar\omega)^{-1}\ln\frac{M+n}{M}
=-(\hbar\omega)^{-1}\ln x$, i.e.\ the very
$\beta=\partial\ln\Omega_{\mathrm{res}}/\partial E$ of
\eqref{eq:expand}--\eqref{eq:beta}, now read off the closed system's own
microcanonical dial, and giving $\langle k\rangle=\mu=(e^{\beta\hbar\omega}-1)^{-1}$,
the Bose occupancy. The two logically distinct steps deserve separating:
\eqref{eq:psw-bound}/\eqref{eq:sho-bound} (typicality, $\rho_S\to\varrho_S$) is
the genuinely new, kinematic content that retires postulate~(i);
\eqref{eq:sho-ratio} ($\varrho_S\to$ Gibbs) is the old large-bath argument---and
it \emph{is} the Taylor truncation of \S\ref{ssec:boltzdist}, since
$\ln\Omega(M{-}k,n{-}1)=(n{-}2)\ln(M{-}k)+\mathrm{const}$ expands to
$-\beta\hbar\omega\,k$ with a dropped quadratic term smaller by $O(1/n)$. The
oscillator lets one \emph{watch} both the postulate and the truncation of
\eqref{eq:expand} become controlled, exact statements.

\paragraph{Why the oscillator is also the cautionary tale.} Typicality is purely
\emph{kinematic}: \eqref{eq:sho-bound} counts states and never used the dynamics.
It guarantees thermal reductions are overwhelmingly abundant on the shell, but
not that an \emph{arbitrary} initial state \emph{evolves into} them---that is the
ergodicity question, and the bare oscillator fails it, instructively, for two
structural reasons. First, within one shell every basis state has the same
energy $\hbar\omega M$, so $e^{-i\hat Ht/\hbar}$ acts as a single global phase
and $\rho_S(t)=\rho_S(0)$ is frozen: the deliberately atypical start of
Remark~\ref{rem:atypical} (all $M$ quanta in oscillator~$1$,
$\rho_S=|M\rangle\langle M|$ pure) stays non-thermal forever. Second, across shells the spectrum is exactly equally spaced, so the
whole system revives exactly every $2\pi/\omega$---there is no dephasing to carry
an out-of-equilibrium state toward the typical set. Equivalently, the eigenstates
are product Fock states $|k_1,\dots,k_n\rangle$, and a single one reduces to a
\emph{pure} $|k_1\rangle$, the opposite of thermal: the harmonic oscillator flatly
violates the eigenstate-thermalisation property a chaotic $\hat H$ would enjoy.
The cure is the one that also makes the spectrum generic---a weak anharmonic
coupling $\lambda\sum_i x_i x_{i+1}$ (the Fermi--Pasta--Ulam--Tsingou
perturbation) lifts the degeneracy and the commensurability, entangles the
eigenstates, and restores both relaxation and eigenstate thermality, at the cost
of the closed-form counting above. The division of labour is exactly the one
worth recording: \emph{typicality} (a theorem, and all the parcel's equilibrium
thermodynamics needs) makes equilibrium states almost all of the shell, while
\emph{ergodicity} (a dynamical hypothesis the integrable oscillator lacks) is the
stronger, separate input needed to \emph{reach} them in time.

\medskip\noindent
None of this changes a line of \texttt{entropy\_S}: the $\beta=1/k_BT$ and the
Boltzmann weights are identical to \S\ref{ssec:boltzdist}. What changes is their
standing---no external reservoir, no imposed postulate, one closed Hamiltonian
system read through a single theorem. This is the same closed-system stance that
\S\ref{sec:conservation} takes for the \emph{conservation} of $s$ on the ascent;
\S\ref{ssec:micro} there and this subsection are its two faces---why the weights
are canonical, and why they stay frozen---each dispensing with the reservoir in
favour of one Hamiltonian.

\subsection{From \texorpdfstring{$Z$}{Z} to entropy, and indistinguishable particles}
\label{ssec:Z}

Everything thermodynamic now follows from $Z$. Insert \eqref{eq:boltzdist} into
the Gibbs entropy \eqref{eq:gibbs}, using $\ln p_i=-\beta E_i-\ln Z$:
\begin{equation}\label{eq:Sderiv}
S=-k_B\sum_i p_i\big(-\beta E_i-\ln Z\big)
 =k_B\beta\sum_i p_iE_i+k_B\ln Z
 =\frac{\langle E\rangle}{T}+k_B\ln Z ,
\end{equation}
where $\langle E\rangle=\sum_i p_iE_i=-\partial\ln Z/\partial\beta$ is the mean
total energy (the thermodynamic internal energy). Defining the Helmholtz free
energy $F\equiv-k_BT\ln Z$, equation \eqref{eq:Sderiv} is
\begin{equation}\label{eq:SfromF}
S=\frac{\langle E\rangle-F}{T}
 =-\Big(\frac{\partial F}{\partial T}\Big)_{V,N},
\end{equation}
the last equality being the standard thermodynamic identity (both $V$ and $N$,
the other natural variables of $F$, are held fixed). So the whole task reduces
to computing $Z$ for the parcel's constituents and differentiating.

\begin{remark}[what depends on what]\label{rem:args}
It is worth being explicit about arguments, as the distinction drives the
conservation argument later. The microstate energies $E_i$ are \emph{mechanical}:
each is set by the Hamiltonian at a fixed volume and particle number, so
$E_i=E_i(V,N)$---a function of the microstate label $i$ and of $(V,N)$, but
\emph{not} of temperature. Consequently the partition function and free energy
carry the mechanical parameters as well, $Z=Z(\beta,V,N)$ and $F=F(T,V,N)$
---which is why $(V,N)$ sit in the derivative subscript of \eqref{eq:SfromF}
---while the \emph{mean} energy $\langle E\rangle=\sum_i p_iE_i=U(T,V,N)$ is a
thermodynamic function of state, belonging to the whole ensemble and to no
single microstate.

A caution on the occupations, since it is easy to over-simplify: $p_i=e^{-\beta
E_i}/Z$ is \emph{not} a function of temperature alone. It depends on $(T,V,N)$
jointly---on $T$ through $\beta$ and on $(V,N)$ through $E_i$---so what fixing
the macrostate determines is the \emph{whole} distribution $\{p_i\}$, with no
residual freedom. This is exactly the structure the reversible-adiabatic
argument of \S\ref{sec:conservation} turns on. Under a slow (quasi-static)
change of $V$ the levels $E_i(V)$ shift, yet along an adiabat the product
$\beta E_i$ is held fixed for \emph{every} $i$ simultaneously---for a monatomic
ideal gas $E_i\propto V^{-2/3}$ and the adiabat is $TV^{2/3}=\mathrm{const}$,
so $\beta E_i\propto T^{-1}V^{-2/3}=\mathrm{const}$---and therefore each
occupation $p_i$, and with it $S=-k_B\sum_i p_i\ln p_i$, is unchanged even as
both $T$ and the $E_i$ move. ``The levels shift; the occupations are frozen'' is
this statement, and it is what makes the ascent isentropic.
\end{remark}

One structural fact makes $Z$ tractable. For $N$ \emph{non-interacting}
particles the total energy is a sum of single-particle energies,
$E_i=\sum_{a=1}^{N}\varepsilon_{a}$, so the system sum factorises into a product
of identical single-particle sums, $z^{N}$. But the particles are
\emph{indistinguishable}: permuting them yields the same physical microstate, so
$z^{N}$ overcounts by the $N!$ permutations. The corrected result---which
resolves the Gibbs mixing paradox and is the classical limit of quantum
statistics---is
\begin{equation}\label{eq:ZN}
Z=\frac{z^{N}}{N!},\qquad
z=\sum_{\text{single-particle states}}e^{-\beta\varepsilon},
\end{equation}
with $z$ the single-particle partition function and $\varepsilon$ a
single-particle energy (contrast $E_i$, the whole-system energy of
\S\ref{ssec:BG}). Everything now hinges on $z$.

\subsection{The bridge to \texorpdfstring{$S(E)$}{S(E)}: Laplace, Legendre, and
why \texorpdfstring{$Z$}{Z} is the elegant object}\label{ssec:bridge}

Two objects describe the same equilibrium: the microcanonical entropy $S(E)$
---geometric and primary, the logarithm of the phase-space volume at fixed
energy---and the canonical $Z(\beta)$ just used. Making the link exact explains
both why they agree and why the $Z$ side carried every derivation above so much
more lightly than the reservoir bookkeeping of \S\ref{ssec:boltzdist} did.

\paragraph{The transform.} Two objects must be kept apart. The \emph{density of
states}
\begin{equation}\label{eq:dos}
\rho(E)=\operatorname{Tr}\,\delta(E-\hat H)=\sum_i\delta(E-E_i)
\end{equation}
of \S\ref{ssec:operator}---with $\delta(E-\hat H)=\sum_i\delta(E-E_i)\,
|i\rangle\langle i|$ the operator function of $\hat H$ (evaluated on the spectrum
and traced \emph{after}, exactly as $Z=\operatorname{Tr}e^{-\beta\hat H}
=\sum_i e^{-\beta E_i}$)---is a distribution, a comb of one spike per level; its
logarithm is meaningless. The Boltzmann entropy instead uses the dimensionless
\emph{count}
\begin{equation}\label{eq:dos-count}
\Omega(E)=\int_E^{E+\Delta E}\!\rho(E')\,\dd E'\approx\bar\rho(E)\,\Delta E,
\qquad S(E)=k_B\ln\Omega(E),
\end{equation}
the number of levels in a macroscopically thin shell $[E,E+\Delta E]$, with
$\bar\rho$ the smoothed envelope of the comb. Because the spectrum is spaced like
$e^{-N}$, any such shell holds astronomically many levels, so $\Omega$ is a huge
\emph{smooth} number and $\ln\Omega$ is well defined and extensive; the window
enters only through the subextensive $k_B\ln\Delta E$, negligible beside
$S\sim k_BN$. ($\Omega$ and $\bar\rho$ thus agree \emph{inside} the logarithm but
are not the same object---they differ by the units-carrying $\Delta E$, which is
why $\rho$ and $\Omega$ get different symbols.) The partition function is the
Laplace transform of the \emph{density} $\rho$,
\begin{equation}\label{eq:bridge-laplace}
Z(\beta)=\int_0^\infty\!\dd E\,\rho(E)\,e^{-\beta E}=\sum_i e^{-\beta E_i}
        \approx\int_0^\infty\!\dd E\,\exp\!\Big[\tfrac{1}{k_B}S(E)-\beta E\Big],
\end{equation}
and the middle equality is the crux: the exponential kernel integrates straight
over the comb, so $Z$ does \emph{for free} the coarse-graining that $S(E)$ needed
a window for---the ``distribution\,$\to$\,analytic function'' elegance, met again
below. It is invertible by the Bromwich integral
$\rho(E)=\frac{1}{2\pi i}\int Z(\beta)\,e^{\beta E}\dd\beta$: $\rho$ and $Z$ carry
identical information. The oscillator shows this concretely---and more cleanly
still, since its spectrum is genuinely discrete, so the ``count'' is an honest
integer degeneracy needing no window at all. The
microcanonical count $\Omega(M,n)=\binom{M+n-1}{n-1}$ of \eqref{eq:sho-dR} is a
genuine combinatorial sequence, but it is packaged whole by
\begin{equation}\label{eq:bridge-genfn}
Z(x)=\sum_{M\ge0}\Omega(M,n)\,x^{M}
     =\sum_{M\ge0}\binom{M+n-1}{n-1}x^{M}
     =\frac{1}{(1-x)^{n}},\qquad x=e^{-\beta\hbar\omega},
\end{equation}
so $Z$ is nothing but the \emph{generating function} whose Taylor coefficients
are the shell dimensions; a single $\Omega(M,n)$ comes back as the coefficient
integral $\frac{1}{2\pi i}\oint Z(x)\,x^{-M-1}\dd x$, whose saddle point gives
the entropy asymptotics. ``Laplace transform of the density of states'' and
``generating function of the counts'' are one statement, continuous and discrete.

\paragraph{The thermodynamic limit is a Legendre transform.} For a macroscopic
system the integrand of \eqref{eq:bridge-laplace} is sharply peaked, and
Laplace's method evaluates it at the saddle $E^\ast(\beta)$ where the exponent is
stationary,
\begin{equation}\label{eq:bridge-saddle}
\frac{\dd}{\dd E}\Big[\tfrac{1}{k_B}S(E)-\beta E\Big]_{E^\ast}=0
\;\Longrightarrow\;
\frac{\dd S}{\dd E}\Big|_{E^\ast}=k_B\beta=\frac1T,
\end{equation}
which is exactly the microcanonical definition of temperature.
Back-substituting, $-k_BT\ln Z=E^\ast-T\,S(E^\ast)=F$: the free energy is the
Legendre transform of $S(E)$, with $\beta$ and $E$ conjugate,
\begin{equation}\label{eq:bridge-legendre}
\ln Z(\beta)=\sup_E\Big[\tfrac{1}{k_B}S(E)-\beta E\Big],\qquad
\tfrac{1}{k_B}S(E)=\inf_\beta\big[\ln Z(\beta)+\beta E\big],
\end{equation}
with duality relations $\beta=\tfrac1{k_B}S'(E)$ and $E=-\partial_\beta\ln Z$.
So the Legendre transform relating the ensembles (the table of
\S\ref{ssec:operator}) is just the $N\to\infty$ shadow of the exact Laplace
transform \eqref{eq:bridge-laplace}: the saddle collapses the whole spectrum onto
one energy $E^\ast$ and one temperature. (This is the thermodynamic face of the
Cram\'er/large-deviation duality between the scaled cumulant generating function
$\ln Z$ and its rate function $S/k_B$.)

\paragraph{Why $Z$ is the elegant object.} Every awkward operation on $S(E)$ is
an easy one on $Z$, because the Laplace transform converts them:
\begin{itemize}
\item \textbf{A distribution becomes an analytic function.}
$\rho(E)=\operatorname{Tr}\delta(E-\hat H)$ is a comb of spectral spikes, and
$S(E)$ needs a coarse-graining window to define at all; $Z(\beta)=\operatorname{Tr}
e^{-\beta\hat H}$ is a smooth, entire function of $\beta$ for any spectrum bounded
below---no delta functions, no window.
\item \textbf{Convolution becomes multiplication---and the reservoir evaporates.}
Independent systems ($\hat H=\hat H_1+\hat H_2$) \emph{convolve} their densities
of states, $\rho=\rho_1\!*\rho_2$, but simply \emph{multiply} their
partition functions, $Z=Z_1Z_2$, so $\ln Z$ is manifestly extensive. The whole
``system\,$+$\,reservoir, expand $\ln\Omega_{\mathrm{res}}(E_{\mathrm{tot}}-E_i)$''
maneuver of \S\ref{ssec:boltzdist} was a convolution evaluated by saddle point;
in the $Z$ picture one never introduces the reservoir---with the intensive
$\beta$ as primary, the Boltzmann weight $e^{-\beta E_i}/Z$ is immediate. The
reservoir's only job was to fix $\beta=S_{\mathrm{res}}'/k_B$, which $Z$ takes as
its argument.
\item \textbf{$\ln Z$ generates the cumulants.} $\langle E\rangle=-\partial_\beta
\ln Z$ and $\operatorname{Var}(E)=\partial_\beta^{2}\ln Z=k_BT^{2}C_V\ge0$; every
fluctuation is a derivative of one smooth function, and fluctuation--dissipation
is one line. Energy fluctuations are not even representable on a single
microcanonical $E$.
\item \textbf{$Z$ is what one computes.} Path integrals, transfer matrices,
perturbation theory, and Monte Carlo all act on $Z$ or $\ln Z$; counting states
at fixed energy is the hard problem. The oscillator is the epitome:
$\Omega(M,n)=\binom{M+n-1}{n-1}$ is a nontrivial count, while $Z=(1-x)^{-n}$ is
one line \eqref{eq:bridge-genfn}.
\end{itemize}

\paragraph{When $S(E)$ is the truer object.} The elegance has a precise price:
the Legendre transform \eqref{eq:bridge-legendre} can only return the
\emph{concave hull} of $S(E)$, so wherever $S(E)$ is non-concave the $Z$ picture
silently discards information---replacing the offending region by its Maxwell
tangent. That is exactly what happens at first-order phase transitions (the flat
tangent is the latent heat), and in long-range or self-gravitating systems, at
negative temperatures, and in small systems---all cases of genuine \emph{ensemble
inequivalence}, where the microcanonical $S(E)$ is faithful and the canonical
average is not. So the honest ranking is the geometric one: $S(E)$ is primary and
information-complete; $Z(\beta)$ is its Laplace transform---smoother, factorising,
fluctuation-generating, and matched to how a system is actually held (at fixed
$T$, not fixed $E$)---and the two are fully equivalent precisely when $S(E)$ is
concave, i.e.\ short-range interactions as $N\to\infty$. That is the parcel's
regime, and the reason the intensive $s(T,p)$---not a fixed-energy count---is the
right variable for the solver.

\paragraph{Hard constraint, soft constraint, and the separation of scales.} The
same transform settles what the microcanonical and canonical descriptions are to
\emph{each other}, and answers a natural worry: the closed whole is governed by
exact energy conservation, yet a small piece of it obeys the Boltzmann
distribution at a fixed temperature---these feel as if they should be the same
statement, and are not quite. Fixing the total energy exactly is a \emph{hard}
constraint---a delta function $\delta(E_{\mathrm{tot}}-\hat H)$ projecting onto
the energy shell, with \emph{zero} energy fluctuation. Fixing only the
\emph{mean} energy through $\beta$ is a \emph{soft} constraint---the Boltzmann
weight $e^{-\beta\hat H}$, with $\beta$ a Lagrange multiplier (the Jaynes reading
of \S\ref{ssec:boltzdist}) and the energy free to fluctuate. They are genuinely
different ensembles, yet locked together by an identity: the hard projector is an
integral of soft weights over imaginary temperature,
\begin{equation}\label{eq:delta-laplace}
\delta(E_{\mathrm{tot}}-\hat H)
=\frac{1}{2\pi i}\int_{c-i\infty}^{c+i\infty}\!\dd\beta\;
   e^{\beta(E_{\mathrm{tot}}-\hat H)},
\end{equation}
which is just the inverse Bromwich transform of \eqref{eq:bridge-laplace} read as
an operator identity---the sharp shell is a coherent superposition of Boltzmann
weights at \emph{every} $\beta$. ``Hard'' and ``soft'' are one function and its
Laplace transform, not two physical laws. What makes them \emph{equivalent} for a
macroscopic system is a separation of scales, entering in two places.

\emph{Whole-system equivalence, guarded by $1/\sqrt N$.} For large $N$ the
$\beta$-integral \eqref{eq:delta-laplace}---equivalently the Laplace integral
\eqref{eq:bridge-laplace}---is dominated by a single real saddle $\beta^\ast$
where $S'(E_{\mathrm{tot}})=1/T$, so the sharp constraint collapses onto
\emph{one} effective temperature. The width of that collapse is the energy
fluctuation the soft constraint permits, which by the cumulant identity
$\operatorname{Var}(E)=\partial_\beta^{2}\ln Z=k_BT^{2}C_V$ above and the
extensivity $C_V,\langle E\rangle\sim N$ is
\begin{equation}\label{eq:relfluct}
\frac{\sqrt{\operatorname{Var}(E)}}{\langle E\rangle}
=\frac{\sqrt{k_BT^{2}C_V}}{\langle E\rangle}\sim\frac{\sqrt N}{N}=\frac{1}{\sqrt N}.
\end{equation}
Energy fluctuations are $O(\sqrt N)$ absolutely but $O(1/\sqrt N)$ relatively: the
soft constraint is \emph{effectively} as sharp as the hard one, and the two
ensembles agree on every intensive quantity to that order.

\emph{Subsystem equivalence, guarded by $N_S/N_B$.} The framing closest to the
parcel is a small subsystem $S$ inside a hard-constrained whole $S+B$.
Marginalising the shell over the bath gives the subsystem's energy distribution
\begin{equation}\label{eq:submarginal}
p(E_S)\;\propto\;\Omega_S(E_S)\,\Omega_B(E_{\mathrm{tot}}-E_S)
        \;=\;\Omega_S(E_S)\,e^{S_B(E_{\mathrm{tot}}-E_S)/k_B},
\end{equation}
which is exactly $p_i\propto\Omega_{\mathrm{res}}(E_{\mathrm{tot}}-E_i)$ of
\eqref{eq:pprop}. Because $E_S\ll E_{\mathrm{tot}}$---that is, $N_S\ll N_B$---the
bath entropy is linear over the range $E_S$ explores,
$S_B(E_{\mathrm{tot}}-E_S)\simeq S_B(E_{\mathrm{tot}})-E_S/T$, and
\eqref{eq:submarginal} becomes the Boltzmann weight
$p(E_S)\propto\Omega_S(E_S)\,e^{-\beta E_S}$: \textbf{the hard constraint on the
whole is a soft constraint on the part.} The bath's temperature is
\emph{stiff}---$\dd\beta/\dd E=-1/(k_BT^{2}C_B)$ with $C_B\sim N_B$---so draining
$E_S$ barely moves $\beta$; the discarded curvature is the finite-bath
back-reaction, of order $E_S/E_{\mathrm{tot}}\sim N_S/N_B$. This linearisation
\emph{is} the Taylor truncation \eqref{eq:expand} performed by hand in
\S\ref{ssec:boltzdist}, now with its small parameter named.

\emph{When the guard fails.} Both equivalences are asymptotic and conditional. If
$S(E)$ is non-concave the saddle of \eqref{eq:delta-laplace} is not unique---two
compete, at a first-order transition---and the ensembles genuinely diverge, the
very ensemble inequivalence just discussed; if the system is small,
\eqref{eq:relfluct} is not negligible and the soft constraint is visibly looser
than the hard one. So the instinct that the two ``ought to be the same thing'' is
right only in the limit: they are one Laplace pair, made \emph{operationally}
interchangeable by the $1/\sqrt N$ and $N_S/N_B$ separations, and distinguishable
again the moment those separations close.

Having set $Z$ against its microcanonical shadow, we return to the object
everything hinges on: the single-particle $z$.

\subsection{The translational partition function and the thermal de~Broglie wavelength}
\label{ssec:lambda}

For a structureless particle in a box of volume $V$ the energy is purely
kinetic, $\varepsilon=\mathbf p^{2}/2m$. In the classical (dilute,
non-degenerate) regime the sum over single-particle states becomes a
phase-space integral with one state per volume $h^{3}$ of phase space:
\begin{equation}\label{eq:ztrans1}
z_{\mathrm{tr}}=\frac{1}{h^{3}}\int_V\!\dd^{3}r\int_{\mathbb R^3}\!\dd^{3}p\;
e^{-\beta \mathbf p^{2}/2m}
=\frac{V}{h^{3}}\prod_{i=x,y,z}\int_{-\infty}^{\infty}\!\dd p_i\,
e^{-\beta p_i^{2}/2m}.
\end{equation}
Each of the three Gaussian integrals evaluates to
$\int_{-\infty}^{\infty}e^{-\beta p^{2}/2m}\dd p=\sqrt{2\pi m/\beta}
=\sqrt{2\pi m k_BT}$, so
\begin{equation}\label{eq:ztrans2}
z_{\mathrm{tr}}=\frac{V\,(2\pi m k_BT)^{3/2}}{h^{3}}\equiv\frac{V}{\Lambda^{3}},
\qquad
\boxed{\;\Lambda\equiv\frac{h}{\sqrt{2\pi m k_BT}}\;}
\end{equation}
The length $\Lambda$ is the \emph{thermal de~Broglie wavelength}: the quantum
coherence scale of a particle with thermal momentum
$\sim\!\sqrt{m k_BT}$. It carries the entire temperature dependence of
$z_{\mathrm{tr}}$, as $\Lambda^{3}\propto T^{-3/2}$. The classical treatment
\eqref{eq:ZN}--\eqref{eq:ztrans1} is valid precisely when $n\Lambda^{3}\ll1$
(mean spacing $\gg$ coherence length); for air at surface conditions
$n\Lambda^{3}\sim10^{-6}$, so the approximation is excellent.

\subsection{Free energy and the Sackur--Tetrode entropy}\label{ssec:sackur}

Take a monatomic gas first, $z=z_{\mathrm{tr}}=V/\Lambda^{3}$ (no internal
structure). From \eqref{eq:ZN}, $\ln Z_N=N\ln z-\ln N!$, and with Stirling's
approximation $\ln N!\simeq N\ln N-N$,
\begin{equation}\label{eq:Fmon}
F=-k_BT\ln Z_N
=-N k_BT\Big[\ln\frac{z}{N}+1\Big]
=-N k_BT\Big[\ln\frac{V}{N\Lambda^{3}}+1\Big].
\end{equation}
The entropy follows from \eqref{eq:SfromF}. Writing
$\ln(V/N\Lambda^{3})=\ln(V/N)+\tfrac32\ln T+C_0$ with
$C_0=\tfrac32\ln(2\pi m k_B/h^{2})$ temperature-independent, only the
$\tfrac32\ln T$ piece and the explicit prefactor $T$ depend on temperature:
\begin{align}
S=-\frac{\partial F}{\partial T}
&=N k_B\Big[\ln\frac{V}{N\Lambda^{3}}+1\Big]
 +N k_BT\cdot\frac{3}{2T} \notag\\
&=N k_B\Big[\ln\frac{V}{N\Lambda^{3}}+\underbrace{1}_{\text{from }F}
      +\underbrace{\tfrac32}_{\partial_T(\tfrac32 T\ln T)}\Big]
 =N k_B\Big[\ln\frac{V}{N\Lambda^{3}}+\frac{5}{2}\Big].
\label{eq:sackur}
\end{align}
This is the \emph{Sackur--Tetrode} equation. Its additive $\tfrac52$ has a
definite provenance: $1$ from the indistinguishability remainder (the $+1$ in
$F$, i.e.\ the $-\ln N!$ Stirling term) and $\tfrac32$ from differentiating the
thermal energy $\tfrac32 N k_BT$. Physically the two variable terms are the
split we will track to the end: $\ln(V/N)$ is \emph{configurational} (positional
microstates---more room per particle, more states), and
$-3\ln\Lambda=\tfrac32\ln T+\mathrm{const}$ is \emph{thermal} (the momentum
distribution broadens with $T$).

\subsection{Intensive form, and the fate of the two \texorpdfstring{$5/2$}{5/2}'s}
\label{ssec:intensive}

Reduce \eqref{eq:sackur} to a specific entropy. Per particle, in units of
$k_B$,
\begin{equation}\label{eq:sperparticle}
\tilde s\equiv\frac{S}{N k_B}=\ln\frac{V}{N\Lambda^{3}}+\frac52 .
\end{equation}
Use the ideal-gas law in the form $V/N=k_BT/p$ and
$\Lambda^{-3}=(2\pi m k_B/h^{2})^{3/2}T^{3/2}$:
\begin{align}
\ln\frac{V}{N\Lambda^{3}}
&=\ln\frac{k_BT}{p}+\frac32\ln T+C_0
 =\underbrace{\ln T}_{\text{from }V/N}\;-\;\ln p\;+\;
   \underbrace{\tfrac32\ln T}_{\text{from }\Lambda^{-3}}\;+\;C_1 \notag\\
&=\frac52\ln T-\ln p+C_1 ,\qquad C_1=\ln k_B+C_0 .
\label{eq:lnVNL}
\end{align}
Substituting into \eqref{eq:sperparticle},
\begin{equation}\label{eq:monatomic}
\tilde s=\frac52\ln T-\ln p+\Big(C_1+\frac52\Big)
        =\underbrace{\frac52}_{=\,c_p/R}\ln T-\ln p+\mathrm{const}.
\end{equation}
There are \emph{two} factors of $5/2$ here and they behave oppositely---this is
the crux:
\begin{itemize}
\item The \textbf{additive} $\tfrac52$ inherited from Sackur--Tetrode
      \eqref{eq:sperparticle} is a pure constant. It is swallowed into the
      reference constant in \eqref{eq:monatomic} and cancels in every entropy
      \emph{difference}. So yes---it drops out of anything physical; it survives
      only in the absolute, third-law-anchored entropy.
\item The \textbf{coefficient} $\tfrac52$ multiplying $\ln T$ is not a constant
      and does \emph{not} drop out. Equation~\eqref{eq:lnVNL} shows it is
      assembled as $1+\tfrac32$: one power of $\ln T$ from the configurational
      $\ln(V/N)=\ln(k_BT/p)$, and $\tfrac32$ from the thermal
      $\Lambda^{-3}\!\propto T^{3/2}$. For a monatomic gas this equals $c_p/R$.
\end{itemize}
Multiplying \eqref{eq:monatomic} by the specific gas constant $R=k_B/m$ (mass
$m$ per particle) turns $\tilde s$ into entropy per unit mass, and
$\tfrac52 R=c_p$ for a monatomic gas, giving the building block
\begin{equation}\label{eq:idealS}
\boxed{\;s \;=\; c_p\ln T - R\ln p + \mathrm{const}.\;}
\end{equation}

\subsection{Beyond monatomic: internal degrees of freedom}\label{ssec:polyatomic}

Sackur--Tetrode as derived assumes a structureless particle, so
\eqref{eq:sackur} does \emph{not} apply verbatim to air ($\mathrm N_2,\mathrm
O_2$) or water vapour ($\mathrm H_2\mathrm O$), which rotate and vibrate. Two
things change, both controlled. First, the single-particle energy is a sum
$\varepsilon=\varepsilon_{\mathrm{tr}}+\varepsilon_{\mathrm{int}}$
(translation plus rotation, vibration, electronic), so the partition function
\emph{factorises},
\begin{equation}\label{eq:zfactor}
z=z_{\mathrm{tr}}(T,V)\,z_{\mathrm{int}}(T),
\end{equation}
with the internal factor a function of $T$ alone (the internal level spacings do
not depend on the box). Second, carrying $z_{\mathrm{int}}$ through
\eqref{eq:Fmon}--\eqref{eq:sackur} adds an internal energy
$u_{\mathrm{int}}=k_BT^{2}\,\dd\ln z_{\mathrm{int}}/\dd T$ per particle and a
matching internal entropy; the extra heat capacity is
$c_{v,\mathrm{int}}=\dd u_{\mathrm{int}}/\dd T$.

The \textbf{equipartition theorem} makes this concrete: every quadratic degree
of freedom that is thermally active contributes $\tfrac12 k_B$ to the
per-particle heat capacity. Translation always gives $3\times\tfrac12$;
rotation adds $2\times\tfrac12$ for a linear molecule ($\mathrm N_2,\mathrm
O_2,\mathrm{CO}_2$) or $3\times\tfrac12$ for a nonlinear one ($\mathrm H_2\mathrm
O$); each active vibrational mode counts twice (kinetic $+$ potential). Hence
$c_v=\tfrac{f}{2}k_B$ over $f$ active quadratic freedoms, $c_p=c_v+k_B$, and the
monatomic ratio $c_p/R=\tfrac52$ is replaced by $c_p/R=\tfrac{f}{2}+1$.

Crucially, the \emph{functional form} of the entropy is unchanged. Whenever
$c_p$ is constant---all active modes classical, none freezing in or out over the
range of interest---integrating the thermodynamic identity
$\dd s=c_p\,\dd T/T-R\,\dd p/p$ (derived independently from the first and second
laws in \S\ref{sec:classical}) gives
\begin{equation}\label{eq:idealSgen}
s=c_p\ln T-R\ln p+\mathrm{const}
\end{equation}
for \emph{any} ideal gas, with the material $c_p$. Sackur--Tetrode is just the
monatomic special case that additionally fixes the absolute constant. For the
atmosphere ($\approx\!200$--$320\K$) the diatomic rotational modes are fully
excited while the stiff vibrational modes are largely frozen, so $\cpd$ and
$\cpv$ (Table~\ref{tab:const}) are very nearly temperature-independent and
\eqref{eq:idealSgen} holds to high accuracy; the small residual $c_p(T)$
variation is absorbed into the effective constants and, for water, into the
$\cl$ tuning of Remark~\ref{rem:modcl}.

\begin{remark}[what the additive constant hides, and why it does not matter]
The reference constant in \eqref{eq:idealS}/\eqref{eq:idealSgen} absorbs
$\Lambda$, the Sackur--Tetrode $\tfrac52$, the internal-mode ground states, and
each species' zero of energy. Because the ascent conserves a \emph{difference}
of entropies at fixed composition (\S\ref{sec:conservation}) and buoyancy
depends on temperature \emph{differences} (\S\ref{sec:adiabat}), these constants
never enter a physical prediction; the code is free to drop them, and does.
\end{remark}

\subsection{The building block, and how it feeds everything below}
\label{ssec:buildingblock}

Equation~\eqref{eq:idealSgen} is the payoff of this section and the answer to
``how does the ideal-gas entropy relate to everything that follows''. It is not
a monatomic curiosity: it is the single functional shape that each gaseous
constituent of the parcel instantiates \emph{once}. Dry air and water vapour are
both ideal gases, so each contributes
\begin{equation}\label{eq:perspecies}
s_k=c_{p,k}\ln T-R_k\ln(\text{its own partial pressure})+\mathrm{const},
\end{equation}
differing only in which $c_p$, which $R$, and---by Dalton's law
(\S\ref{sec:parcel})---which partial pressure it sees. The total specific
entropy of the parcel is then the mixing-ratio-weighted sum of these identical
blocks plus a condensed-phase term (Gibbs additivity, \S\ref{sec:parcel}); once
phase equilibrium is used to reintroduce the latent heat (\S\ref{sec:phase}),
that sum is exactly Emanuel's \eqref{eq:E94}. In short: derive the block once
here, and ``everything below'' is bookkeeping---which partial pressure, which
$c_p$, which mixing ratio---layered on top of \eqref{eq:idealSgen}.

\section{Classical scaffolding: why \texorpdfstring{$s$}{s} is a state function}
\label{sec:classical}

Statistical mechanics gives $s$; classical thermodynamics tells us it is a
\emph{function of state}---indispensable, since the solver evaluates $s$ at a
point $(T,p,\rvv)$ with no reference to how the parcel got there.

\subsection{First and second laws}

The first law for a closed system is $\dd U=\delta Q-\delta W$, with
$\delta W=p\,\dd V$ for quasi-static volume work. Neither $\delta Q$ nor
$\delta W$ is an exact differential: $\oint \delta Q\neq 0$. The second law
supplies an \emph{integrating factor}. In Carath\'eodory's formulation, the
empirical fact that arbitrarily close to any equilibrium state there exist
states not reachable by any adiabatic path ($\delta Q=0$) is equivalent, for
the Pfaffian form $\delta Q$, to the existence of an integrating denominator;
that denominator is the absolute temperature $T$, and
\begin{equation}\label{eq:dS}
\dd s \;=\; \frac{\delta Q_{\mathrm{rev}}}{T}
\end{equation}
is an \emph{exact} differential. Its integral, $s$, is a state function. This
is the precise sense in which ``entropy exists''.

\subsection{Recovering \texorpdfstring{$s=c_p\ln T-R\ln p$}{s=cp lnT - R lnp} thermodynamically}

For an ideal gas, $pV=nRT$ and $U=U(T)$ with $\dd U=c_v\,\dd T$ (Joule's law).
Then $\delta Q_{\mathrm{rev}}=\dd U+p\,\dd V=c_v\,\dd T+p\,\dd V$, and with
$p/T=R/v$ (specific volume $v$),
\begin{equation}
\dd s=\frac{c_v\,\dd T}{T}+\frac{R\,\dd v}{v}
     =\frac{c_p\,\dd T}{T}-\frac{R\,\dd p}{p},
\end{equation}
using Mayer's relation $c_p-c_v=R$ and $\dd(\ln pv)=\dd\ln(RT)$ to switch
variables. Integrating returns \eqref{eq:idealS}, now with the material $c_p$.
Statistical mechanics and classical thermodynamics agree, as they must; the
exactness of $\dd s$ (equivalently the Maxwell relation
$\partial^2 s/\partial T\partial p=\partial^2 s/\partial p\partial T$) is what
lets \texttt{entropy\_S} be a pure function of its three arguments.

\section{The parcel as a multicomponent mixture}\label{sec:parcel}

\subsection{Composition and the per-dry-air convention}

The parcel is dry air plus water in vapour and (once saturated) condensed form.
Masses are bookkept as \emph{mixing ratios} relative to the dry-air mass $m_d$:
\begin{equation}
\rvv=\frac{m_{\text{vapour}}}{m_d},\qquad
\rl=\frac{m_{\text{liquid}}}{m_d},\qquad
\rt=\rvv+\rl .
\end{equation}
Entropy is quoted \emph{per unit mass of dry air} because $m_d$ is the one mass
that is conserved along the ascent (water changes phase, but dry air is inert
and neither created nor removed); dividing by it gives a bookkeeping base that
does not move. This is why \eqref{eq:target} carries $\Rd$ bare but weights the
water terms by mixing ratios.

\subsection{Dalton's law and the vapour--pressure relation}

Each gas in an ideal mixture behaves as if it alone occupied the volume, so the
total pressure is the sum of partial pressures (Dalton): $p=\pd+e$, with $e$
the vapour partial pressure and $\pd=p-e$ the dry-air partial pressure. The
mixing ratio and the partial pressure are two encodings of the same amount of
vapour; equating $\rvv=m_v/m_d=(\eps^{-1}e)/(p-e)$ with $\eps=\Rd/\Rv$ gives the
mutually inverse maps
\begin{equation}\label{eq:evrv}
e=\frac{\rvv\,p}{\eps+\rvv},\qquad
\rvv=\frac{\eps\,e}{p-e},
\end{equation}
which are exactly \texttt{utilities.ev} and \texttt{utilities.rv}. The
appearance of $p-e$ (not $p$) inside the dry-air logarithm of \eqref{eq:target}
is Dalton's law at work: the dry air sees only its own partial pressure.

\subsection{Additivity of entropy: Gibbs' theorem}

For an ideal-gas mixture the entropy is the sum of the entropies each component
would have alone at its own partial pressure and the common temperature
(Gibbs' theorem---there is no entropy of mixing beyond what the partial
pressures already encode, because ideal gases do not interact). Hence we may
build $s$ additively:
\begin{equation}\label{eq:additive}
s \;=\; s_d \;+\; \rvv\, s_v \;+\; \rl\, \sigma_\ell,
\end{equation}
$s_d$ the specific entropy of dry air, $s_v$ that of vapour, and $\sigma_\ell$
that of liquid water, each per unit mass of its own species. We now need the
three pieces.

\section{Water substance and phase equilibrium}\label{sec:phase}

The gaseous pieces use \eqref{eq:idealS}:
\begin{equation}\label{eq:sd-sv}
s_d = \cpd\ln T - \Rd\ln \pd + \mathrm{const}_d,\qquad
s_v = \cpv\ln T - \Rv\ln e + \mathrm{const}_v.
\end{equation}
Liquid water is nearly incompressible, so its entropy is essentially thermal,
\begin{equation}\label{eq:sl}
\sigma_\ell = \cl\ln T + \mathrm{const}_\ell .
\end{equation}
What ties the vapour and liquid constants together---and produces the latent
term in \eqref{eq:target}---is phase equilibrium.

\subsection{Chemical potential and the Clausius--Clapeyron relation}

At a flat interface, vapour and liquid coexist when their specific Gibbs free
energies (chemical potentials per unit mass) are equal, $g_v(T,\es)=g_\ell(T)$,
which \emph{defines} the saturation vapour pressure $\es(T)$. Differentiating
along the coexistence curve with $\dd g=-s\,\dd T+v\,\dd p$ gives
\begin{equation}
-s_v\,\dd T + v_v\,\dd\es \;=\; -\sigma_\ell\,\dd T + v_\ell\,\dd\es
\;\Longrightarrow\;
\frac{\dd\es}{\dd T}=\frac{s_v-\sigma_\ell}{v_v-v_\ell}.
\end{equation}
The reversible latent heat is the enthalpy of the isothermal, isobaric phase
change, $\Lv=T(s_v-\sigma_\ell)$, i.e.\ the \emph{entropy of vaporisation} is
\begin{equation}\label{eq:svap}
s_v-\sigma_\ell=\frac{\Lv}{T}.
\end{equation}
With $v_v\gg v_\ell$ and $v_v=\Rv T/\es$ (ideal vapour), this is the
Clausius--Clapeyron equation
\begin{equation}\label{eq:cc}
\frac{\dd\es}{\dd T}=\frac{\Lv\,\es}{\Rv\,T^{2}}.
\end{equation}

\subsection{Kirchhoff's relation and the coded fits}

The latent heat itself depends on temperature. Applying
$\dd\Lv/\dd T=\cpv-\cl$ (Kirchhoff's law: the two phases warm at different
specific heats) and integrating from $0^{\circ}$C,
\begin{equation}\label{eq:kirchhoff}
\Lv(T)=\Lv^{0}+(\cpv-\cl)\,(T-T_0),
\end{equation}
which is exactly \texttt{utilities.Lv} (Table~\ref{tab:const}:
$\Lv^0=2.501\times10^6$, slope $\cpv-\cl=-630$). Substituting
\eqref{eq:kirchhoff} into \eqref{eq:cc} and integrating yields a closed
expression for $\es(T)$; the code uses the algebraically cheaper
August--Roche--Magnus fit
\begin{equation}\label{eq:esfit}
\es(T)=6.112\,\exp\!\Big(\frac{17.67\,T_C}{243.5+T_C}\Big)\ \mathrm{hPa},
\qquad T_C=T-273.15,
\end{equation}
i.e.\ \texttt{utilities.es\_cc}, which reproduces the Clausius--Clapeyron
integral to sub-percent accuracy over atmospheric temperatures. Equations
\eqref{eq:svap}--\eqref{eq:esfit} are all we need from phase physics.

\begin{remark}[the ``modified'' $\cl=2500\JkgK$]\label{rem:modcl}
The true specific heat of liquid water is $\approx4190\JkgK$; \texttt{tcpyPI}
(following \texttt{pcmin}) uses $2500$. This is a deliberate effective value
that folds the ice phase and the temperature dependence of $\Lv$ into a single
linear Kirchhoff slope $\cpv-\cl=-630$, tuned so that \eqref{eq:kirchhoff}
tracks the observed $\Lv(T)$ across the mixed-phase range the parcel traverses.
It changes the reference thermodynamics slightly but not the structure of any
derivation here.
\end{remark}

\section{Assembling the total specific entropy}\label{sec:assembly}

Insert \eqref{eq:sd-sv}--\eqref{eq:sl} into the additive decomposition
\eqref{eq:additive} and use $\rl=\rt-\rvv$:
\begin{align}
s &= \big[\cpd\ln T-\Rd\ln\pd\big]
   + \rvv\big[\cpv\ln T-\Rv\ln e\big]
   + (\rt-\rvv)\,\cl\ln T + \mathrm{const} \notag\\
  &= \cpd\ln T-\Rd\ln\pd
   + \rt\,\cl\ln T
   + \rvv\big[(\cpv-\cl)\ln T-\Rv\ln e\big] + \mathrm{const}.
\label{eq:collected}
\end{align}
The bracket on the vapour term is where the latent heat re-enters. Evaluate
\eqref{eq:svap} using \eqref{eq:sd-sv}--\eqref{eq:sl} \emph{at saturation}
($e=\es$):
\begin{equation}\label{eq:latentconst}
\frac{\Lv}{T}=s_v-\sigma_\ell
=(\cpv-\cl)\ln T-\Rv\ln\es+(\mathrm{const}_v-\mathrm{const}_\ell).
\end{equation}
Solve \eqref{eq:latentconst} for $(\cpv-\cl)\ln T$ and substitute into the
bracket of \eqref{eq:collected}:
\begin{align}
(\cpv-\cl)\ln T-\Rv\ln e
&=\Big[\tfrac{\Lv}{T}+\Rv\ln\es-(\mathrm{const}_v-\mathrm{const}_\ell)\Big]-\Rv\ln e
\notag\\
&=\frac{\Lv}{T}-\Rv\ln\frac{e}{\es}-(\mathrm{const}_v-\mathrm{const}_\ell)
 =\frac{\Lv}{T}-\Rv\ln H+\mathrm{const}.
\end{align}
Putting this back into \eqref{eq:collected} and absorbing all reference
constants,
\begin{equation}\label{eq:E94}
\boxed{\;s=(\cpd+\rt\,\cl)\ln T-\Rd\ln(p-e)+\frac{\Lv\,\rvv}{T}-\rvv\,\Rv\ln H.\;}
\end{equation}
This is Emanuel (1994), eq.~4.5.9, and term-for-term it is \eqref{eq:target}
and \texttt{utilities.entropy\_S}. The four terms carry the four physical
meanings we tracked from the start:
\begin{center}
\begin{tabular}{ll}
\toprule
Term & Physical content \\
\midrule
$(\cpd+\rt\cl)\ln T$ & thermal entropy of dry air $+$ all condensate carried\\
$-\Rd\ln(p-e)$ & configurational entropy of the dry air (Dalton partial pressure)\\
$+\Lv\rvv/T$ & latent entropy stored in the vapour fraction (Clausius--Clapeyron)\\
$-\rvv\Rv\ln H$ & dilution entropy of sub-saturated vapour ($H<1\Rightarrow{>}0$)\\
\bottomrule
\end{tabular}
\end{center}

\begin{remark}[the single-argument \texttt{R}, and reduction below saturation]
The code's \eqref{eq:target} uses one mixing ratio \texttt{R} in every water
term, i.e.\ it identifies $\rvv=\rt=R$. That is exact wherever the parcel is
unsaturated, which is precisely where \texttt{entropy\_S} is called: at the
launch state, to set the conserved reference $s_0$. Below the lifting
condensation level all water is vapour, $\rl=0$, and \eqref{eq:E94} is the
familiar dry/unsaturated entropy with $H<1$. The generalisation to $\rvv<\rt$
(condensate present) is carried implicitly by the saturated form used in the
inversion, \eqref{eq:SGcode} below, where $H=1$ retires the last term.
\end{remark}

\section{Why \texorpdfstring{$s$}{s} is conserved on the ascent}
\label{sec:conservation}

\subsection{The exact entropy budget}

For the parcel as a closed system, the Clausius inequality splits any change of
entropy into a reversible flux and an irreversible production:
\begin{equation}\label{eq:budget}
\dd s \;=\; \underbrace{\frac{\delta Q}{T}}_{\text{flux}}
        \;+\; \underbrace{\dd s_{\mathrm{irr}}}_{\ge 0\ \text{(production)}} .
\end{equation}
The parcel model annihilates both terms by two idealisations.

\paragraph{Adiabatic: no flux, $\delta Q=0$.} Convective ascent is fast
compared with the time for heat to diffuse or radiate across the parcel
boundary. Formally, the ratio of the advective timescale $H_\rho/w$ (density
scale height over vertical velocity, minutes) to the thermal-diffusion and
radiative timescales (hours to days) is $\ll1$, so $\delta Q\approx0$ over the
ascent. The first term of \eqref{eq:budget} vanishes.

\paragraph{Reversible: no production, $\dd s_{\mathrm{irr}}=0$.} The parcel is
assumed to remain in internal thermodynamic equilibrium---uniform $T$ and $p$,
no turbulent entrainment or mixing with the environment, and, crucially, phase
change occurring exactly at saturation. At saturation the chemical potentials
are equal, $g_v=g_\ell$ (\S\ref{sec:phase}), so transferring water across the
phase boundary produces \emph{no} entropy: condensation onto a parcel held at
$H=1$ is reversible. (Condensing sub- or super-saturated vapour would carry a
finite $\Delta g$ per unit mass and generate $\dd s_{\mathrm{irr}}>0$.) The
second term of \eqref{eq:budget} vanishes.

Together, $\delta Q=0$ and $\dd s_{\mathrm{irr}}=0$ give
\begin{equation}\label{eq:isentropic}
\boxed{\;\dd s = 0 \quad\text{along the ascent}\;}
\end{equation}
---the ascent is \emph{isentropic}. This is the entire justification for
holding \texttt{entropy\_S} fixed.

\subsection{Closed for water: reversible vs.\ pseudo-adiabatic}

Conservation of \eqref{eq:E94} as a function of $(T,p)$ alone additionally
requires that the water inventory not change: the ``reversible'' adiabat keeps
all condensate suspended in the parcel, so $\rt=\mathrm{const}$ and the parcel
is a closed system for water. In the code this is the default
\texttt{ascent\_flag}$=0$; the density temperature (\S\ref{sec:adiabat}) then
carries the full water loading $\rt$. The alternative
\texttt{ascent\_flag}$=1$ (\emph{pseudo}-adiabatic) removes condensate the
instant it forms: the system is then \emph{open}, $\rt$ is not conserved, and
the invariant is not \eqref{eq:E94} but an approximate pseudo-entropy. The
distinction is physical, not numerical.

\subsection{The microscopic reading}\label{ssec:micro}

Statistical mechanics gives \eqref{eq:isentropic} a clean interpretation. A
reversible adiabatic change is a \emph{slow} deformation of the parcel's
Hamiltonian (volume and level spacings change as $p$ and $T$ change) with no
coupling to a heat bath. Two theorems then apply. By \textbf{Liouville's
theorem}, phase-space volume is invariant under Hamiltonian flow; by the
\textbf{adiabatic theorem}, a sufficiently slow deformation preserves the
occupation probability of each state (no transitions between energy shells).
The accessible phase-space volume $\Omega$ is therefore unchanged, and by
\eqref{eq:boltzmann} so is $S$; equivalently every $p_i$ in the Gibbs form
\eqref{eq:gibbs} is frozen, so $-k_B\sum p_i\ln p_i$ is constant even as the
individual energy levels move. Isentropic literally means \emph{the
microstate count does not change---the levels merely shift}. Note that the two
senses of ``adiabatic'' (thermodynamic: no heat; mechanical: slow) here
coincide, which is why the slow, heat-free ascent lands on a single adiabat.
This is the closed-system twin of \S\ref{ssec:typicality}: there the same
no-reservoir, single-Hamiltonian stance fixes \emph{which} weights $p_i$ are
canonical (a pure state's reduced density matrix); here it explains why those
weights are \emph{frozen} under the reversible ascent. Together they carry both
halves of the Gibbs entropy \eqref{eq:gibbs}---the value of the weights and
their constancy---without ever invoking a heat bath.

\section{The moist adiabat and buoyancy}\label{sec:adiabat}

\subsection{From $\dd s=0$ to a path}

Because $s$ is a state function, \eqref{eq:isentropic} is a differential
constraint that pins the path in the $(T,p)$ plane:
\begin{equation}\label{eq:pathODE}
0=\dd s=\frac{\partial s}{\partial T}\,\dd T+\frac{\partial s}{\partial p}\,\dd p
\;\Longrightarrow\;
\frac{\dd T}{\dd p}=-\frac{\partial s/\partial p}{\partial s/\partial T}.
\end{equation}
This ODE is the moist adiabat. Two regimes:

\paragraph{Unsaturated (below the LCL): the dry adiabat.} With $\rvv=\rt$ fixed
and no condensation, the dominant balance of \eqref{eq:E94} is
$s\approx(\cpd+\rt\cl)\ln T-\Rd\ln\pd+\mathrm{const}$. Holding $s$ and the water
constant and treating the small vapour corrections as constant, $\dd s=0$ gives
$c_p\,\dd\ln T=\Rd\,\dd\ln p$ (with the appropriate moist $c_p$), i.e.\ Poisson's
relation
\begin{equation}\label{eq:poisson}
T\propto p^{\,\Rd/c_p},\qquad
\theta\equiv T\Big(\frac{p_0}{p}\Big)^{\Rd/c_p}=\mathrm{const},
\end{equation}
the conserved (virtual) potential temperature. This is exactly the branch the
CAPE routine takes below the condensation level, where it writes
$T_G=T_P\,(p/p_P)^{\Rd/\cpd}$.

\paragraph{The lifting condensation level.} As the unsaturated parcel cools
along \eqref{eq:poisson}, $\es(T)$ falls (Clausius--Clapeyron) while $e$ falls
only with pressure; $H=e/\es$ rises to $1$. The pressure at which $H=1$ is the
\emph{lifting condensation level} (LCL), \texttt{utilities.e\_pLCL} in the
code. Above it the parcel is saturated.

\paragraph{Saturated (above the LCL): the moist adiabat proper.} Now
$\rvv=r_s(T,p)=\eps\es(T)/(p-\es)$ is pinned to saturation and
$\rl=\rt-r_s$ is condensate. Substituting into \eqref{eq:pathODE} (equivalently
differentiating \eqref{eq:E94} at $H=1$) gives the saturated adiabatic lapse
rate; its steepness is set by the latent-heat release, and the temperature
sensitivity is dominated by $\dd r_s/\dd T$ through Clausius--Clapeyron. Rather
than integrate this ODE, the solver conserves $s$ algebraically---the inversion
of \S\ref{sec:inversion}.

\subsection{Density temperature and buoyancy}

The dynamically relevant temperature is the \emph{density temperature}
$\Trho$, defined so that $p=\rho\,\Rd\,\Trho$ holds with the dry-air gas
constant despite the moisture. Accounting for vapour (which lightens the air,
$\Rv>\Rd$) and suspended condensate (which weighs it down),
\begin{equation}\label{eq:trho}
\Trho=T\,\frac{1+\rvv/\eps}{1+\rt},
\end{equation}
exactly \texttt{utilities.Trho(T,\,$\rt$,\,$\rvv$)}. The buoyancy that drives
the parcel---and whose vertical integral is the CAPE of the companion
document---is the density-temperature excess over the environment,
\begin{equation}\label{eq:buoy}
b=\Rd\,\big(\Trho^{\text{parcel}}-\Trho^{\text{env}}\big)\big/\,\text{(a length)},
\end{equation}
i.e.\ the parcel's self-determined adiabat (\S\ref{sec:conservation}) compared
against the sounding. The environment enters only here, in the comparison; the
parcel's own $T(p)$ came entirely from conserving $s$.

\section{The inversion: recovering \texorpdfstring{$T$}{T} from \texorpdfstring{$s$}{s}}
\label{sec:inversion}

\subsection{Definition and well-posedness}

At each new pressure level the solver knows the conserved entropy $s_0$ and
must find the temperature consistent with it. This is a function
\emph{inversion}: given the forward map $T\mapsto s(T,p)$ at fixed $p$ (and
$\rt$), solve
\begin{equation}\label{eq:invert}
s(T,p)=s_0 \qquad\text{for } T=s^{-1}(s_0;p).
\end{equation}
An inverse exists and is unique wherever the forward map is strictly monotone.
Differentiate the saturated entropy---the relevant branch for the inversion,
\begin{equation}\label{eq:SGcode}
s_{\mathrm{sat}}(T)=(\cpd+\rt\,\cl)\ln T-\Rd\ln(p-\es)+\frac{\Lv\,r_s}{T},
\end{equation}
(the $-\rvv\Rv\ln H$ term is gone because $H=1$)---with respect to $T$ at fixed
$p$. Using Clausius--Clapeyron \eqref{eq:cc} for $\dd r_s/\dd T$, the dominant
terms give
\begin{equation}\label{eq:dsdT}
\frac{\partial s_{\mathrm{sat}}}{\partial T}
=\frac{1}{T}\left(\cpd+\rt\,\cl+\frac{\Lv^{2}\,r_s}{\Rv\,T^{2}}\right) \;>\;0.
\end{equation}
Every term is manifestly positive, so $s_{\mathrm{sat}}(\cdot)$ is strictly
increasing and \eqref{eq:invert} has a unique root. By the inverse function
theorem the inverse is itself $C^1$ with derivative
$(s^{-1})'=1/(\partial s/\partial T)$. The inversion is well posed---this is the
``smooth, strictly monotone entropy relation'' invoked without proof in the
companion analysis.

\subsection{Newton iteration, and its appearance in the code}

There is no closed form for $s^{-1}$, so the solver inverts numerically by
Newton--Raphson on the residual $F(T)=s_{\mathrm{sat}}(T)-s_0$:
\begin{equation}\label{eq:newton}
T_{n+1}=T_n-\frac{F(T_n)}{F'(T_n)}
       =T_n+\frac{s_0-s_{\mathrm{sat}}(T_n)}{\partial s_{\mathrm{sat}}/\partial T}.
\end{equation}
This is exactly \texttt{solve\_temperature\_from\_entropy}: its variable
\texttt{SG} is $s_{\mathrm{sat}}(T_n)$ of \eqref{eq:SGcode}, its \texttt{SL} is
the derivative \eqref{eq:dsdT},
\[
\texttt{SL}=\frac{\cpd+\rt\,\cl+\Lv^2 r_s/(\Rv T^2)}{T},
\]
and the update \texttt{TGNEW = TG + AP*(S-SG)/SL} is \eqref{eq:newton} with a
relaxation factor \texttt{AP} (smaller for the first steps, then unity) guarding
against overshoot from a poor initial guess. Quadratic convergence follows from
$F\in C^2$ and $F'\neq0$, both guaranteed by \eqref{eq:dsdT}. The positivity of
\texttt{SL} is not a numerical convenience; it is the inverse-function-theorem
hypothesis, and \eqref{eq:dsdT} shows it holds term by term.

\begin{remark}[why this is the ``entropy relation'', not a lapse-rate integral]
One could instead integrate the moist-adiabat ODE \eqref{eq:pathODE} in $p$.
Conserving $s$ algebraically via \eqref{eq:invert} is preferable: it carries no
accumulated integration error, it makes the conserved quantity exact at every
level by construction, and it exposes the monotonicity \eqref{eq:dsdT} that
guarantees a unique parcel temperature. The price---a root-find per level---is
cheap and quadratically convergent.
\end{remark}

\section{Synthesis}\label{sec:synthesis}

The chain is now complete and first-principled at every link:
\begin{enumerate}
\item Entropy is a microstate count, $S=k_B\ln\Omega=-k_B\sum p_i\ln p_i$
      (\S\ref{sec:statmech}); for an ideal gas it evaluates to
      $s=c_p\ln T-R\ln p+\mathrm{const}$ (Sackur--Tetrode).
\item The second law makes $s$ an exact state function via the integrating
      factor $1/T$ (\S\ref{sec:classical}), so it may be evaluated pointwise.
\item The parcel is a Dalton mixture; entropy adds over components at their
      partial pressures (\S\ref{sec:parcel}), and phase equilibrium
      ($g_v=g_\ell$) yields Clausius--Clapeyron and the vaporisation entropy
      $\Lv/T$ (\S\ref{sec:phase}).
\item Summing dry air, vapour, and condensate gives Emanuel's total specific
      entropy \eqref{eq:E94}---\texttt{utilities.entropy\_S} exactly
      (\S\ref{sec:assembly}).
\item It is conserved because the ascent is adiabatic (no flux) and reversible
      (no production, phase change at saturation), and reversible ascent keeps
      $\rt$ fixed (\S\ref{sec:conservation}); microscopically, Liouville plus
      the adiabatic theorem freeze $\Omega$.
\item Conserving $s$ defines the moist adiabat \eqref{eq:pathODE} and, via the
      strictly monotone \eqref{eq:dsdT}, a well-posed inversion
      $s(T,p)=s_0\mapsto T$ that \texttt{solve\_temperature\_from\_entropy}
      performs by guarded Newton iteration (\S\ref{sec:adiabat}--\ref{sec:inversion}).
\end{enumerate}
Every constant, fit, and function named above appears in the \texttt{tcpyPI}
source---in \texttt{constants.py}, \texttt{utilities.py}, and \texttt{pi.py};
the physics that justifies each is derived here rather than asserted. The one modelling choice
worth flagging for a physicist reading the code is the effective
$\cl=2500\JkgK$ (Remark~\ref{rem:modcl}); everything else is textbook
thermodynamics assembled in the meteorologist's per-dry-air bookkeeping.

\end{document}
